\documentclass[
  12pt,
  a4paper,
  brazil,
  twoside
]{report}

\usepackage[utf8]{inputenc}
\usepackage[T1]{fontenc}
\usepackage{times}
\usepackage[brazil]{babel}

\usepackage[
  top=3cm, left=3cm, right=2cm, bottom=2cm
]{geometry}

\usepackage{setspace}
\onehalfspacing

\usepackage{indentfirst}
\setlength{\parindent}{1.25cm}

\usepackage{fancyhdr}
\pagestyle{fancy}
\fancyhf{}
\fancyhead[LE,RO]{\thepage}
\fancyhead[RE,LO]{\textit{A Insustent\''avel Leveza da Detec\c{c}\~ao}}
\renewcommand{\headrulewidth}{0.5pt}

\usepackage[colorlinks=true, linkcolor=black, citecolor=black, urlcolor=black]{hyperref}
\usepackage{array, booktabs, longtable}
\usepackage{graphicx}
\usepackage[bottom]{footmisc}

\usepackage{titlesec}
\titleformat{\chapter}{\normalfont\bfseries\Large}{\thechapter}{0.5em}{}
\titleformat{\section}{\normalfont\bfseries\large}{\thesection}{0.5em}{}
\titleformat{\subsection}{\normalfont\bfseries\normalsize}{\thesubsection}{0.5em}{}

\usepackage{tocloft}
\renewcommand{\cftchapleader}{\cftdotfill{\cftdotsep}}
\renewcommand{\cftsecleader}{\cftdotfill{\cftdotsep}}
\setcounter{tocdepth}{2}

\begin{document}

\begin{titlepage}
  \begin{center}
    \vspace*{2cm}
    {\large\bfseries UNIVERSIDADE DE S\~AO PAULO\par}
    {\large\bfseries PROGRAMA DE P\'OS-GRADUA\c{C}\~AO EM CI\^ENCIA DA INFORMA\c{C}\~AO\par}
    \vspace{4cm}
    {\LARGE\bfseries A Insustentável Leveza da Detecção: Regulação da
Inteligência Artificial Generativa na Pesquisa Científica Brasileira
entre a Norma, a Técnica e o Mercado Global da Quarta Revolução
Industrial\par}
    \vspace{0.5cm}
    {\large\bfseries Análise Crítica da Portaria CNPq nº 2.664/2026\par}
    \vspace{3cm}
    {\large Disserta\c{c}\~ao apresentada ao Programa de P\'os-Gradua\c{c}\~ao em
    Ci\^encia da Informa\c{c}\~ao da Universidade de S\~ao Paulo como requisito
    parcial para obten\c{c}\~ao do t\'itulo de Mestre.\par}
    \vspace{1cm}
    {\large Orientador: Prof. Dr. [Nome]\par}
    \vfill
    {\large 2026\par}
  \end{center}
\end{titlepage}

\newpage
\begin{center}
  \vspace*{4cm}
  {\Large\bfseries FOLHA DE ROSTO\par}
  \vspace{2cm}
  {\normalsize Disserta\c{c}\~ao de Mestrado\par}
  \vspace{0.5cm}
  {\normalsize \textbf{A Insustentável Leveza da Detecção: Regulação da
Inteligência Artificial Generativa na Pesquisa Científica Brasileira
entre a Norma, a Técnica e o Mercado Global da Quarta Revolução
Industrial}\par}
  \vspace{0.5cm}
  {\normalsize \textbf{Análise Crítica da Portaria CNPq nº
2.664/2026}\par}
  \vspace{1cm}
  \begin{flushright}
  \begin{minipage}{8cm}
  \small Disserta\c{c}\~ao apresentada ao Programa de P\'os-Gradua\c{c}\~ao em
  Ci\^encia da Informa\c{c}\~ao da Universidade de S\~ao Paulo como requisito
  parcial para obten\c{c}\~ao do t\'itulo de Mestre.
  \end{minipage}
  \end{flushright}
  \vfill
  {\large 2026\par}
\end{center}

\newpage
\begin{center}
  \vspace*{2cm}
  {\Large\bfseries RESUMO\par}
  \vspace{1cm}
  \begin{minipage}{14cm}
  \onehalfspacing
  Esta dissertação analisa criticamente a Portaria CNPq nº 2.664, de 6
  de março de 2026, que institui a Política de Integridade na Atividade
  Científica do CNPq, à luz da Quarta Revolução Industrial. A pesquisa
  adota metodologia jurídico-dogmática de natureza qualitativa,
  combinando análise normativa, revisão sistemática da literatura
  técnica sobre detectores de Inteligência Artificial Generativa (IAG),
  regulação comparada (União Europeia, China, Estados Unidos) e análise
  de economia política da inovação. Os resultados revelam sete
  ambiguidades hermenêuticas no art. 9º da Portaria, inviabilidade
  técnica da fiscalização baseada em detectores automáticos, assimetria
  competitiva no mercado global de P\&D e conflitos normativos com a
  LGPD e a Constituição Federal. Conclui-se que a Portaria, não obstante
  suas boas intenções, impõe custos regulatórios desproporcionais ao
  sistema brasileiro de CT\&I e propõe-se modelo híbrido de regulação
  que combine marcação técnica pelos provedores de IA, declaração
  simplificada pelo pesquisador e gradação de exigências por nível de
  risco.
  \end{minipage}
  \vspace{1cm}
  \par\noindent
  \begin{minipage}{14cm}
  \textbf{Palavras-chave:} Inteligência Artificial Generativa;
Integridade Científica; CNPq; Quarta Revolução Industrial; Regulação
Comparada.
  \end{minipage}
\end{center}

\newpage
\begin{center}
  \vspace*{2cm}
  {\Large\bfseries ABSTRACT\par}
  \vspace{1cm}
  \begin{minipage}{14cm}
  \onehalfspacing
  This thesis critically analyzes CNPq Ordinance No.~2.664, of March 6,
  2026, which establishes the Integrity Policy for Scientific Activity,
  in light of the Fourth Industrial Revolution. The research adopts a
  qualitative legal-dogmatic methodology, combining normative analysis,
  systematic review of the technical literature on Generative AI
  detectors, comparative regulation (European Union, China, United
  States), and political economy analysis of innovation. The results
  reveal seven hermeneutic ambiguities in Article 9 of the Ordinance,
  technical infeasibility of enforcement based on automatic detectors,
  competitive asymmetry in the global RD market, and normative conflicts
  with the Brazilian Data Protection Law and the Federal Constitution.
  The conclusion is that the Ordinance, despite its good intentions,
  imposes disproportionate regulatory costs on the Brazilian STI system.
  A hybrid regulation model is proposed, combining technical
  watermarking by AI providers, simplified researcher declaration, and
  risk-based requirements.
  \end{minipage}
  \vspace{1cm}
  \par\noindent
  \begin{minipage}{14cm}
  \textbf{Keywords:} Generative Artificial Intelligence; Scientific
Integrity; CNPq; Fourth Industrial Revolution; Comparative Regulation.
  \end{minipage}
\end{center}

\newpage
\tableofcontents
\newpage

\pagenumbering{arabic}

\section{1. Introdução}\label{introduuxe7uxe3o}

O ano de 2026 será lembrado como o momento em que três das maiores
potências científicas do mundo --- Brasil, China e União Europeia ---
convergiram para regular o uso de inteligência artificial generativa
(IAG) na pesquisa acadêmica, cada qual com seu modelo normativo. Os
Estados Unidos, por sua vez, optaram por via distinta, centrada em
segurança de pesquisa e combate à influência estrangeira. Esta
dissertação situa a Portaria CNPq nº 2.664/2026 neste cenário comparado
para revelar suas fragilidades, limites e propor alternativa regulatória
viável.

\subsection{1.1. Contexto Regulatório}\label{contexto-regulatuxf3rio}

A Portaria CNPq nº 2.664, de 6 de março de 2026, institui a Política de
Integridade na Atividade Científica. Seu art. 9º estabelece três
diretrizes centrais: (i) alínea ``c'': obrigatoriedade de declarar o uso
de IAG ``em qualquer fase do desenvolvimento da pesquisa (concepção,
redação, análise de dados, submissão) especificando {[}\ldots{]} a
ferramenta utilizada e a finalidade''; (ii) alínea ``d'': vedação de
``submissão de conteúdo gerado por IAG como se fosse de autoria
humana''; (iii) alínea ``e'': vedação ao uso de IAG para elaboração de
pareceres científicos \footnote{BRASIL. Portaria CNPq nº 2.664, de 6 de
  março de 2026. Institui a Política de Integridade na Atividade
  Científica. \emph{Diário Oficial da União}, Brasília, 11 mar. 2026. A
  Portaria reconhece oficialmente a IAG como ferramenta legítima, não a
  proíbe, alinhando-se ao movimento internacional. Sua fragilidade
  central é não definir sanções específicas nem mecanismos de
  verificação, criando norma sem enforcement --- lex imperfecta. A
  ausência de tipificação sancionatória viola o princípio da legalidade
  estrita (art. 5º, II, CF) ao expor o pesquisador ao arbítrio
  interpretativo do agente público.}. Simultaneamente, a CAPES
disponibilizou o documento ``A inteligência artificial na pesquisa e no
fomento: desafios e oportunidades'' (Brasil, 2026) \footnote{BRASIL, A.
  L. \emph{A inteligência artificial na pesquisa e no fomento: desafios
  e oportunidades}. Brasília: CAPES, 2026. Documento interno de
  circulação restrita. Disponível mediante solicitação à CAPES. Sua
  principal contribuição é sinalizar que a agência reconhece a IA como
  tema prioritário.}, que reconhece a urgência do debate mas não oferece
solução operacional.

A pergunta fundamental --- ``como provar que alguém usou IA?'' --- não
encontra resposta no texto da Portaria nem na literatura técnica. Esta
dissertação sustenta que, ao exigir o que é logicamente impossível de
verificar, o protocolo CNPq opera como instrumento de retardamento
burocrático da pesquisa, penalizando inovadores e não constrangendo
mal-intencionados.

\subsection{1.2. Problema de Pesquisa}\label{problema-de-pesquisa}

A Portaria CNPq nº 2.664/2026 pressupõe que o uso não declarado de IAG
pode ser identificado e punido. Esta premissa é questionável à luz de
três ordens de problemas: (a) a inviabilidade técnica dos detectores de
texto gerado por IA; (b) as ambiguidades hermenêuticas do texto
normativo; (c) a ausência de mecanismos de enforcement compatíveis com o
estado da arte da tecnologia. A estas se soma a constatação de que o
modelo regulatório brasileiro difere substancialmente das abordagens
adotadas por China, União Europeia e Estados Unidos, criando assimetria
competitiva no mercado global de pesquisa e desenvolvimento.

\subsection{1.3. Objetivos}\label{objetivos}

O objetivo geral desta dissertação é analisar criticamente a Portaria
CNPq nº 2.664/2026 à luz da literatura técnica sobre detecção de IAG, da
teoria jurídica e da regulação comparada. São objetivos específicos: (a)
sistematizar as evidências empíricas sobre a confiabilidade dos
detectores automatizados de IAG; (b) identificar e categorizar as
ambiguidades hermenêuticas do art. 9º; (c) analisar as colisões
normativas com a LGPD e a Constituição Federal; (d) comparar o modelo
brasileiro com as regulações da China, União Europeia e Estados Unidos;
(e) propor modelo regulatório híbrido que supere as limitações
identificadas.

\subsection{1.4. Metodologia}\label{metodologia}

A pesquisa adota metodologia jurídico-dogmática de natureza qualitativa,
combinando: (a) análise normativa da Portaria CNPq nº 2.664/2026 e
legislação correlata; (b) revisão sistemática da literatura técnica
publicada entre 2023 e 2026 sobre detectores de IAG, com ênfase em
estudos empíricos de larga escala; (c) análise de regulação comparada
compreendendo China (GB 45438-2025, Interim Measures), União Europeia
(AI Act Regulation 2024/1689) e Estados Unidos (NSF TRUST Framework);
(d) análise de economia política da inovação à luz do conceito de Quarta
Revolução Industrial (Schwab, 2016).

\subsection{1.5. Estrutura da
Dissertação}\label{estrutura-da-dissertauxe7uxe3o}

Além desta introdução, a dissertação compreende cinco capítulos
adicionais. O Capítulo 2 analisa a dimensão técnica do problema da
detecção de IAG. O Capítulo 3 examina as ambiguidades hermenêuticas do
art. 9º e os paradoxos éticos da declaração obrigatória. O Capítulo 4
investiga as colisões normativas com o ordenamento jurídico brasileiro.
O Capítulo 5 apresenta a análise comparada dos quatro modelos
regulatórios. O Capítulo 6 conclui com balanço crítico e propostas de
aperfeiçoamento.

\section{2. Dimensão Técnica: O Problema Insolúvel da
Detecção}\label{dimensuxe3o-tuxe9cnica-o-problema-insoluxfavel-da-detecuxe7uxe3o}

\subsection{2.1. A Falência dos Detectores
Automatizados}\label{a-faluxeancia-dos-detectores-automatizados}

O art. 9º, alínea ``c'', pressupõe que o uso não declarado de IAG pode
ser identificado e punido. Esta premissa é falseada por um corpo robusto
de evidências empíricas.

Weber-Wulff et al.~(2023) testaram 14 detectores comerciais de texto
gerado por IA e encontraram ``altas taxas de falso-positivo'',
concluindo que ``nenhum detector é confiável o suficiente para uso em
contextos de alta penalidade'' \footnote{WEBER-WULFF, D. et al.~Testing
  of detection tools for AI-generated text. \emph{International Journal
  for Educational Integrity}, v. 19, n.~1, 2023. DOI:
  10.1007/s40979-023-00146-z. Estudo financiado pelo BMBF (Alemanha).
  Limitação: testou apenas detectores disponíveis em 2023; versões
  posteriores podem ter desempenho diferente, mas a literatura de
  2024-2025 confirma que o problema persiste. Metodologia
  multiplataforma robusta com 14 ferramentas.}. A amostra incluiu textos
acadêmicos reais e gerados por diferentes modelos (GPT-3, GPT-3.5,
GPT-4), em múltiplos idiomas, e os resultados foram consistentes:
nenhuma ferramenta atingiu simultaneamente sensibilidade superior a 80\%
e especificidade superior a 90\%.

Liu et al.~(2024) compararam detectores automatizados com julgamento
humano na identificação de textos médicos gerados por LLMs. O estudo
concluiu que ``nenhum método atinge confiabilidade aceitável'' e que ``a
concordância entre detectores é baixa'' \footnote{LIU, J. Q. J. et
  al.~The great detectives: humans versus AI detectors in catching large
  language model-generated medical writing. \emph{International Journal
  for Educational Integrity}, v. 20, n.~1, 2024. DOI:
  10.1007/s40979-024-00155-6. Originalidade: compara detecção humana
  versus automatizada em contexto médico. Limitação: amostra pequena
  (n=200 textos).}.

Pudasaini et al.~(2025), em survey abrangente no Journal of Academic
Ethics (Qualis A1), documentaram que a paráfrase --- técnica trivial
disponível em qualquer LLM --- reduz a acurácia dos detectores a níveis
próximos do aleatório. Os autores testaram cinco detectores contra
textos submetidos a paráfrase, tradução e prompt injection, constatando
que ``a taxa de detecção cai de 85\% para 32\% após paráfrase mínima''
\footnote{PUDASAINI, S. et al.~Survey on AI-generated plagiarism
  detection. \emph{Journal of Academic Ethics}, v. 23, p.~1137-1170,
  2025. DOI: 10.1007/s10805-024-09576-x. Qualis A1. Mapeia
  sistematicamente 47 detectores, 23 datasets e 18 estratégias de
  evasão. Limitação: foco em inglês; idiomas como português têm
  desempenho de detecção ainda pior.}.

O estudo de larga escala mais recente é o de Liang et al.~(2025), que
analisou 164.579 artigos completos e 5,2 milhões de resumos. A conclusão
é devastadora para o protocolo CNPq: entre os 75.172 artigos publicados
desde 2023, apenas 76 (\textasciitilde0,1\%) declararam uso de IA.
Contudo, a proporção estimada de artigos com conteúdo gerado por IA no
primeiro trimestre de 2025 era 40 vezes maior que a taxa de declaração
(40:1) \footnote{LIANG, W. et al.~Discrepancy between AI content
  proportion and disclosure rate in scientific publications.
  \emph{Nature Human Behaviour}, 2025. DOI: 10.48550/arXiv.2512.06705. A
  escala do estudo (5,2 milhões de artigos) confere poder estatístico
  ímpar. A razão 40:1 é o dado mais contundente contra a eficácia do
  protocolo brasileiro.}.

Kwon (2025), em reportagem na Nature, documentou que ``science sleuths''
identificaram centenas de artigos científicos publicados em 2024-2025
que usam IAG sem declarar, corroborando a discrepância identificada por
Liang et al. \footnote{KWON, D. Science sleuths flag hundreds of papers
  that use AI without disclosing it. \emph{Nature}, v. 641, p.~290-291,
  2025. DOI: 10.1038/d41586-025-01180-2.}.

\subsection{2.2. A Anomalia do Art. 9º, Alínea
``c''}\label{a-anomalia-do-art.-9uxba-aluxednea-c}

A redação da alínea ``c'' --- ``IAG, de qualquer espécie e em qualquer
fase'' --- padece de hipergeneralização. A IA não é monolítica. Kocak et
al.~(2025) distinguem ao menos quatro níveis de uso de IA na pesquisa
acadêmica, cada qual com diferentes implicações éticas e de detecção
\footnote{KOCAK, B. et al.~Ensuring peer review integrity in the era of
  large language models. \emph{European Journal of Radiology Artificial
  Intelligence}, v. 2, 100018, 2025. DOI: 10.1016/j.ejrai.2025.100018.
  Artigo com 47 co-autores de 12 países. A taxonomia de níveis de uso de
  IA proposta pelos autores (Seção 2) é referência internacional.}:

\begin{enumerate}
\def\labelenumi{\arabic{enumi}.}
\tightlist
\item
  Assistência linguística (corretor ortográfico, paráfrase pontual): não
  detectável, relevância ética mínima.
\item
  Sumarização e tradução (DeepL, Elicit): parcialmente detectável,
  relevância ética baixa.
\item
  Geração de seções (redação de métodos ou discussão): detecção
  duvidosa, relevância ética alta.
\item
  Geração integral (artigo completo por IA): detectável apenas por
  marcadores grosseiros, relevância ética gravíssima.
\end{enumerate}

A Portaria trata todos os níveis como equivalentes. Esta falta de
gradação expõe o pesquisador que usa Grammarly para correção ortográfica
ao mesmo regime legal de quem submete artigo integralmente gerado por IA
--- um tratamento juridicamente desproporcional e que viola o princípio
da razoabilidade.

\subsection{2.3. A Falácia do
Enforcement}\label{a-faluxe1cia-do-enforcement}

A impossibilidade técnica de detecção produz uma falácia de enforcement:
o protocolo parece exigir conduta, mas não oferece meio de verificação.
Pérez-Neri et al.~(2025), na Frontiers in Artificial Intelligence,
sintetizam: ``os desenvolvimentos atuais {[}em detectores{]} não são
suficientemente precisos ou confiáveis para uso em avaliação acadêmica''
\footnote{PÉREZ-NERI, I. et al.~Detecting and correcting errors in
  scientific literature in the generative AI era. \emph{Frontiers in
  Artificial Intelligence}, v. 8, 2025. DOI: 10.3389/frai.2025.1644098.
  Revisão de escopo com análise de 22 ferramentas de detecção.}.

Gils et al.~(2024) demonstraram, em estudo de policy prototyping para o
AI Act europeu, que mesmo requisitos de transparência bem definidos
encontram barreiras técnicas de implementação, desde a falta de
padronização até a impossibilidade de watermarking eficaz em modelos de
linguagem de pequeno porte \footnote{GILS, T. et al.~From Policy to
  Practice: Prototyping The EU AI Act's Transparency Requirements.
  \emph{SSRN}, 2024. DOI: 10.2139/ssrn.4714345.}.

\section{3. Dimensão Hermenêutica e
Ética}\label{dimensuxe3o-hermenuxeautica-e-uxe9tica}

\subsection{3.1. As Sete Ambiguidades do Art.
9º}\label{as-sete-ambiguidades-do-art.-9uxba}

O art. 9º da Portaria nº 2.664/2026 contém pelo menos sete ambiguidades
que comprometem sua aplicação coerente e produzem insegurança jurídica:

\textbf{Ambiguidade 1 --- ``pesquisa apoiada pelo CNPq'' (caput).} A
expressão pode ser interpretada restritivamente (apenas pesquisas com
repasse financeiro direto) ou ampliativamente (toda atividade de
pesquisador que recebe bolsa PQ). A primeira interpretação exclui
parcela significativa da produção científica brasileira; a segunda
amplia o escopo além da competência legal do CNPq, invadindo seara de
periódicos e universidades \footnote{PESQUISAI. CNPq reconhece o uso da
  Inteligência Artificial Generativa na pesquisa científica.
  \emph{Portal PesquisAI}, 16 mar. 2026.}.

\textbf{Ambiguidade 2 --- ``qualquer espécie'' de IAG (alínea ``c'').}
Inclui corretor ortográfico baseado em IA? Tradutor automático? A falta
de definição captura desde ferramentas assistivas triviais até modelos
generativos de fronteira, sem qualquer gradação.

\textbf{Ambiguidade 3 --- ``qualquer fase'' (alínea ``c'').} A concepção
da ideia de pesquisa já dispara a obrigação de declarar? A redação
sugere que sim, mas a operacionalização é inviável.

\textbf{Ambiguidade 4 --- ``conteúdo gerado por IAG'' (alínea ``d'').}
Esta é a ambiguidade mais problemática. Se o pesquisador usa ChatGPT
para sugerir três versões de um parágrafo, edita manualmente as três e
funde em uma quarta versão, o texto resultante é ``gerado por IAG''? A
literatura sobre human-AI co-creation demonstra que a distinção é um
continuum, não uma dicotomia.

\textbf{Ambiguidade 5 --- ``finalidade'' a especificar (alínea ``c'').}
Basta ``revisão textual'' ou exige detalhamento por parágrafo? A
Portaria não estabelece o nível de granularidade exigido.

\textbf{Ambiguidade 6 --- ``ferramenta utilizada'' (alínea ``c'').}
Versão exata do modelo? Provedor? Data da consulta? Parâmetros
utilizados (temperatura, top-p)?

\textbf{Ambiguidade 7 --- vedação a pareceres com IAG (alínea ``e'').}
Aplica-se a LLMs locais (open-source) executados sem transmissão de
dados? A ascensão de modelos como Llama 3.1 405B, DeepSeek-V3 e Mistral
Large viabiliza execução local com desempenho comparável a modelos
comerciais, tornando a vedação inócua para usuários técnicos
\footnote{A ascensão de modelos open-source (Llama 3.1 405B,
  DeepSeek-V3, Mistral Large) viabiliza execução local com desempenho
  comparável a modelos comerciais, tornando a alínea ``e'' inócua para
  usuários técnicos.}.

\subsection{3.2. O Paradoxo da
Transparência}\label{o-paradoxo-da-transparuxeancia}

Schilke e Reimann (2025) identificaram o que denominam ``dilema da
transparência'' (transparency dilemma): a declaração de uso de IA reduz
a confiança do leitor no trabalho, mesmo quando a IA foi usada de forma
ética e assistiva \footnote{SCHILKE, O.; REIMANN, M. The transparency
  dilemma: How AI disclosure erodes trust. \emph{Organizational Behavior
  and Human Decision Processes}, v. 188, 2025. DOI:
  10.1016/j.obhdp.2025.104405. Estudo experimental com 6 estudos (N
  total = 2.847).}. O estudo, conduzido com 2.847 participantes em seis
experimentos, demonstrou que o efeito persiste mesmo quando a IA é usada
apenas para formatação ou revisão gramatical.

Este efeito é particularmente perverso no contexto do protocolo CNPq: o
pesquisador que declara uso de IA segundo o art. 9º, alínea ``c'',
expõe-se a escrutínio mais rigoroso, vieses de revisão e potencial
desqualificação. O pesquisador que não declara, por outro lado, não
corre risco (a detecção é inviável). O resultado é a seleção adversa: o
protocolo pune a honestidade e recompensa a omissão.

\subsection{3.3. O Paradoxo da Reprodutibilidade versus Segredo
Industrial}\label{o-paradoxo-da-reprodutibilidade-versus-segredo-industrial}

A alínea ``d'' --- vedação de submissão de conteúdo gerado por IAG como
autoria humana --- colide com a realidade da pesquisa de fronteira
mediada por IA. Modelos de AlphaFold, GNoME ou descoberta por IA operam
sob lógica de caixa-preta protegida por segredo industrial.

Kocak et al.~(2025) descrevem este paradoxo: a reprodutibilidade
científica exige transparência total, mas os modelos algorítmicos que
viabilizam inovações aceleradas são protegidos como propriedade
intelectual. A comunidade pode validar resultados finais, não o processo
cognitivo da máquina \footnote{KOCAK, B. et al., op. cit.}. A Portaria,
ao tratar a transparência como sempre possível, ignora o conflito
estrutural entre propriedade corporativa e escrutínio acadêmico.

Meszaros, Huys e Ioannidis (2026) demonstram que este problema é global:
a distinção entre ``pesquisa'' e ``atividade comercial'' no AI Act
europeu é igualmente problemática, e as isenções de pesquisa
``baseiam-se em distinções que podem não capturar plenamente as
realidades da pesquisa contemporânea em IA'' \footnote{MESZAROS, J.;
  HUYS, I.; IOANNIDIS, J. P. A. Challenges in applying the EU AI Act
  research exemptions to contemporary AI research. \emph{npj Digital
  Medicine}, 2026. DOI: 10.1038/s41746-025-02263-0.}.

\section{4. Dimensão Jurídico-Brasileira: Colisões
Normativas}\label{dimensuxe3o-juruxeddico-brasileira-colisuxf5es-normativas}

\subsection{4.1. LGPD versus Alínea
``c''}\label{lgpd-versus-aluxednea-c}

A alínea ``c'' exige que o pesquisador especifique ``a ferramenta
utilizada e a finalidade'' do uso de IAG. Se a ferramenta for um LLM
comercial (ChatGPT, Claude, Gemini), o pesquisador pode estar --- sem
saber --- violando a Lei Geral de Proteção de Dados Pessoais (LGPD, Lei
nº 13.709/2018) ao submeter dados de pesquisa a servidores no exterior,
especialmente se os dados envolverem informações pessoais de sujeitos de
pesquisa \footnote{A LGPD (Lei nº 13.709/2018) exige consentimento
  específico e finalidade determinada para tratamento de dados pessoais.
  A submissão de dados a LLMs comerciais (que tipicamente usam dados
  para treinamento) pode configurar violação. Cf. art. 7º, I c/c art.
  11, I da LGPD.}.

A LGPD exige consentimento específico e finalidade determinada para
tratamento de dados pessoais (art. 7º, I c/c art. 11, I). A submissão de
dados a LLMs comerciais --- que tipicamente utilizam dados para
treinamento --- pode configurar violação. Adicionalmente, a decisão do
Supremo Tribunal Federal no RE 1.057.258 (Tema 987) reforça a
autodeterminação informativa como direito fundamental.

A Portaria oferece zero orientação sobre como conciliar a obrigação de
declarar detalhadamente o uso de IA com a proteção de dados. O
pesquisador fica diante de um conflito normativo: cumprir o CNPq
(declarar a ferramenta) ou cumprir a LGPD (não expor dados a terceiros
não autorizados pelo consentimento).

\subsection{4.2. Art. 218 da Constituição
Federal}\label{art.-218-da-constituiuxe7uxe3o-federal}

O art. 218 da Constituição Federal determina que ``o Estado promoverá e
incentivará o desenvolvimento científico, a pesquisa, a capacitação
científica e tecnológica e a inovação''. A Portaria impõe obrigações
declaratórias e vedações sem contrapartida de incentivo.

Bucci (2021) sustenta que a regulação estatal da pesquisa deve observar
a proporcionalidade entre meios restritivos e fins promocionais
\footnote{BUCCI, M. P. D. \emph{Fundamentos para uma teoria jurídica das
  políticas públicas}. 2. ed.~São Paulo: Saraiva Jur, 2021.}. Uma norma
que cria custo burocrático sem benefício demonstrável de integridade é
desproporcional na modalidade ``proibição de excesso''. O descompasso
entre regulação e fomento pode ser questionado à luz do princípio
constitucional.

\subsection{4.3. Ausência de Sanção Específica e Legalidade
Estrita}\label{ausuxeancia-de-sanuxe7uxe3o-especuxedfica-e-legalidade-estrita}

A Portaria não estabelece sanções para o descumprimento do art. 9º. Esta
lacuna viola o princípio da legalidade estrita (art. 5º, II, CF) e a
tipicidade no direito administrativo sancionador. Sem descrição clara de
infração e sanção, a aplicação da norma dependeria de interpretação
extensiva --- vedada neste âmbito.

O direito administrativo brasileiro exige que a sanção esteja prevista
em lei ou ato normativo equivalente (art. 5º, XXXIX, CF c/c art. 1º da
LINDB). Justen Filho (2024) esclarece que não há sanção sem prévia
cominação legal (nulla poena sine lege) \footnote{JUSTEN FILHO, M.
  \emph{Curso de direito administrativo}. 15. ed.~Rio de Janeiro:
  Forense, 2024.}. A Portaria CNPq é ato administrativo normativo, não
lei em sentido formal. Mesmo que se admita poder disciplinar da
autarquia (art. 37, §6º, CF), a ausência de especificação do que
constitui infração e qual a penalidade viola o devido processo legal
substantivo.

A ausência de sanção explícita significa que, na prática, a Portaria é
lex imperfecta: obriga mas não pune, ou pune sem tipificação, o que é
juridicamente insustentável.

\subsection{4.4. O Sistema Qualis-CAPES como Antecedente
Estrutural}\label{o-sistema-qualis-capes-como-antecedente-estrutural}

A Portaria nº 2.664/2026 não surge em vácuo. Ela se insere em tradição
de avaliação quantitativa e burocrática que sistematicamente privilegia
métricas internacionais em detrimento da pesquisa nacional. Caballero
Rivero, Santos e Trzesniak (2024) demonstraram que os sistemas de
avaliação CAPES/CNPq induzem padrões de publicação que priorizam
periódicos estrangeiros de alto impacto em detrimento da relevância
local \footnote{CABALLERO RIVERO, A.; SANTOS, R. N. M.; TRZESNIAK, P.
  Efeitos dos sistemas de avaliação de pesquisa de CAPES e CNPQ nos
  padrões de publicação dos pesquisadores das ciências da saúde no
  Brasil. \emph{Em Questão}, v. 30, 2024. DOI:
  10.1590/1808-5245.30.138437.}.

O efeito performativo é bem descrito por Albuquerque (2025): ``o simples
fato de existirem critérios externos faz com que os programas de
pós-graduação se reorganizem para atendê-los'' \footnote{ALBUQUERQUE, U.
  P. Ciência em debate: o peso silencioso da avaliação. \emph{The
  Conversation}, 13 out. 2025.}. Sguissardi (2006) denominou este
fenômeno de ``avaliação defensiva'' \footnote{SGUISSARDI, V. A avaliação
  defensiva no ``modelo CAPES de avaliação''. \emph{Perspectiva}, v. 24,
  n.~1, p.~49-88, 2006.}. Perez (2025) atualizou a crítica para o
contexto do Qualis 2025-2028, argumentando que os novos critérios
``tendem a aprofundar desigualdades regionais, institucionais e de
gênero'' \footnote{PEREZ, O. C. Avaliação em Disputa: A Reforma do
  Qualis e os Desafios para a Ciência. \emph{Novos Debates}, v. 11,
  n.~1, 2025. DOI: 10.48006/2358-0097/V11N1.E111011.}.

\section{5. Regulação Comparada: Quatro Modelos
Regulatórios}\label{regulauxe7uxe3o-comparada-quatro-modelos-regulatuxf3rios}

A comparação internacional revela que cada país abordou o desafio de
forma diversa, e que a opção brasileira --- autodeclaração sem
verificação técnica --- é a de menor densidade normativa e menor
capacidade de enforcement.

\subsection{5.1. China: Regulação Centralizada com Padrão
Obrigatório}\label{china-regulauxe7uxe3o-centralizada-com-padruxe3o-obrigatuxf3rio}

A China adotou o modelo mais abrangente e tecnicamente sofisticado. O
Ministério da Ciência e Tecnologia publicou em 2023 diretrizes que
proíbem o uso de IAG para gerar materiais de submissão, exigem rotulagem
clara de conteúdo gerado por IA e vedam a IA como co-autora \footnote{CHINA.
  Ministry of Science and Technology. \emph{Guidelines for Responsible
  Research Conduct}. Beijing, 2023. Publicado em 21/12/2023. Exige que
  conteúdo gerado por IA seja claramente rotulado. Limitação:
  diretrizes, não lei; enforcement depende de incorporação por
  instituições.}. A Academia Chinesa de Ciências complementou com oito
``lembretes de integridade'' detalhando usos permitidos e proibidos por
etapa de pesquisa.

O marco mais significativo é o padrão nacional obrigatório GB 45438-2025
(``Cybersecurity Technology: Labeling Method for Content Generated by
Artificial Intelligence''), vigente desde setembro de 2025, que exige
marcação explícita (visível ao usuário) e implícita (metadados,
watermarks digitais) de conteúdo gerado por IA \footnote{CHINA.
  \emph{Measures for the Identification of AI-Generated (Synthetic)
  Content}. Promulgado conjuntamente por CAC, MIIT, MPS e NRTA em
  07/03/2025, vigente desde 01/09/2025, acompanhado do padrão nacional
  GB 45438-2025.}. A este se somam as Interim Measures on Generative AI
Services, que estabelecem regime de registro e filing de algoritmos.

O modelo chinês é o que mais se aproxima de um enforcement factível, por
três razões: (a) o padrão técnico obrigatório é aplicado na fonte pelos
provedores de IA chineses; (b) a fiscalização é exercida pelo Cyberspace
Administration of China (CAC) com poder de sanção administrativa; (c) o
regime exige registro e filing de algoritmos de IA.

Limitações do modelo chinês incluem: (a) dependência de provedores
nacionais --- LLMs estrangeiros não são acessíveis na China continental,
criando ecossistema controlado; (b) alto custo de conformidade; (c)
potencial para censura e controle político disfarçado de regulação de
IA, como observa Beladiya (2026) \footnote{BELADIYA, R. How China Is
  Controlling Generative AI Technologies. \emph{AI Law Guide}, maio
  2026.}.

\subsection{5.2. União Europeia: Abordagem Baseada em Risco com
Enforcement
Escalonado}\label{uniuxe3o-europeia-abordagem-baseada-em-risco-com-enforcement-escalonado}

O EU AI Act (Regulation 2024/1689) é o primeiro marco regulatório
abrangente de IA no mundo. Para a pesquisa acadêmica, seus dispositivos
mais relevantes são: o art. 2(6) e (8), que estabelece isenções para
pesquisa (research exemption); o art. 50, que impõe obrigações de
transparência aplicáveis a todos os sistemas de IA; e o Code of Practice
on Marking and Labelling, em desenvolvimento \footnote{EUROPEAN UNION.
  Regulation (EU) 2024/1689. \emph{Official Journal of the European
  Union}, 2024. Art. 2(6): ``This Regulation does not apply to AI
  systems specifically developed and put into service for the sole
  purpose of scientific research and development.'' Art. 50(2):
  ``Providers of GPAI models that generate synthetic content shall
  ensure that the output is marked in a machine-readable format and
  detectable as artificially generated or manipulated.''}.

A distinção crítica do modelo europeu é que não impõe declaração pelo
autor, mas sim marcação técnica pelo provedor. O LLM comercial é
obrigado por lei (art. 50(2)) a marcar seu output de forma legível por
máquina. A obrigação de transparência recai sobre quem disponibiliza o
sistema, não sobre quem o usa.

Isso contrasta fundamentalmente com o modelo brasileiro, que impõe a
obrigação exclusivamente sobre o pesquisador-usuário --- sem exigência
correspondente sobre os provedores. A assimetria é evidente: o Brasil
regula quem não tem como verificar a própria declaração.

Meszaros, Huys e Ioannidis (2026) identificaram fragilidades nas
isenções de pesquisa do AI Act: a distinção entre pesquisa acadêmica e
atividade comercial é turva, e o desenvolvimento em ambiente controlado
versus teste em ambiente real é difícil de determinar \footnote{MESZAROS,
  J.; HUYS, I.; IOANNIDIS, J. P. A., op. cit.}.

\subsection{5.3. Estados Unidos: Segurança de Pesquisa em vez de
Regulação de
IA}\label{estados-unidos-seguranuxe7a-de-pesquisa-em-vez-de-regulauxe7uxe3o-de-ia}

Os EUA não possuem regulação federal específica sobre uso de IA na
pesquisa acadêmica. A abordagem é indireta, centrada em pesquisa
segurança (research security) e combate à influência estrangeira. A NSF
Scientific Integrity Policy (NSF 24-007, fevereiro de 2024) reafirma
compromisso com integridade mas não menciona IA especificamente
\footnote{NSF. \emph{NSF Scientific Integrity Policy} (NSF 24-007). 12
  fev. 2024. Política genérica que estabelece expectativas para manter e
  aprimorar a Integridade Científica na NSF. Silêncio sobre IA reflete a
  posição do governo Biden (EO 14110 de outubro de 2023).}. O NSPM-33 e
o TRUST Framework focam em disclosure de apoio estrangeiro e conflitos
de compromisso \footnote{NSF. \emph{Trusted Research Using Safeguards
  and Transparency (TRUST)}. Jun.~2024. Framework de árvore decisória
  para avaliar riscos de segurança nacional. Nenhum branch trata de IA.}.

O modelo americano é o mais liberal: não exige declaração de IA, não
estabelece padrão de marcação, não cria obrigação para provedores. A
integridade é delegada às instituições e periódicos, que têm políticas
próprias. Esta abordagem minimiza custos burocráticos mas transfere o
problema para o self-governance de cada entidade.

\subsection{5.4. Tabela Comparativa}\label{tabela-comparativa}

{\def\LTcaptype{none} % do not increment counter
\begin{longtable}[]{@{}
  >{\raggedright\arraybackslash}p{(\linewidth - 8\tabcolsep) * \real{0.1827}}
  >{\raggedright\arraybackslash}p{(\linewidth - 8\tabcolsep) * \real{0.1777}}
  >{\raggedright\arraybackslash}p{(\linewidth - 8\tabcolsep) * \real{0.2183}}
  >{\raggedright\arraybackslash}p{(\linewidth - 8\tabcolsep) * \real{0.2234}}
  >{\raggedright\arraybackslash}p{(\linewidth - 8\tabcolsep) * \real{0.1827}}@{}}
\toprule\noalign{}
\begin{minipage}[b]{\linewidth}\raggedright
Dimensão
\end{minipage} & \begin{minipage}[b]{\linewidth}\raggedright
Brasil (CNPq 2.664/2026)
\end{minipage} & \begin{minipage}[b]{\linewidth}\raggedright
China
\end{minipage} &
\multicolumn{2}{>{\raggedright\arraybackslash}p{(\linewidth - 8\tabcolsep) * \real{0.4061} + 2\tabcolsep}@{}}{%
\begin{minipage}[b]{\linewidth}\raggedright
União Europeia \textbar{} EUA (NSF) \textbar{}
\end{minipage}} \\
\midrule\noalign{}
\endhead
\bottomrule\noalign{}
\endlastfoot
Base legal &
\multicolumn{2}{>{\raggedright\arraybackslash}p{(\linewidth - 8\tabcolsep) * \real{0.3959} + 2\tabcolsep}}{%
Portaria administrativa \textbar{} Lei + Padrão Nacional} &
\multicolumn{2}{>{\raggedright\arraybackslash}p{(\linewidth - 8\tabcolsep) * \real{0.4061} + 2\tabcolsep}@{}}{%
Regulamento (diretamente aplicável) \textbar{} Política interna NSF
\textbar{}} \\
Obrigação & Autodeclaração do pesquisador &
\multicolumn{3}{>{\raggedright\arraybackslash}p{(\linewidth - 8\tabcolsep) * \real{0.6244} + 4\tabcolsep}@{}}{%
Marcação pelo provedor + declaração \textbar{} Marcação técnica pelo
provedor (Art. 50)\textbar{} Nenhuma específica \textbar{}} \\
Padrão técnico & Nenhum &
\multicolumn{3}{>{\raggedright\arraybackslash}p{(\linewidth - 8\tabcolsep) * \real{0.6244} + 4\tabcolsep}@{}}{%
GB 45438-2025 (obrigatório) \textbar{} Code of Practice (voluntário)
\textbar{} Nenhum \textbar{}} \\
Enforcement & Inexistente (sem sanção) & CAC com sanções administrativas
&
\multicolumn{2}{>{\raggedright\arraybackslash}p{(\linewidth - 8\tabcolsep) * \real{0.4061} + 2\tabcolsep}@{}}{%
Comissão + autoridades nacionais \textbar{} Research security (não IA)
\textbar{}} \\
Escopo & Pesquisa apoiada pelo CNPq & Toda pesquisa na China &
\multicolumn{2}{>{\raggedright\arraybackslash}p{(\linewidth - 8\tabcolsep) * \real{0.4061} + 2\tabcolsep}@{}}{%
Sistemas de IA no mercado EU \textbar{} Pesquisas financiadas NSF} \\
Custo de conformidade & Alto (declaração detalhada) & Médio (rotulagem +
verificação) &
\multicolumn{2}{>{\raggedright\arraybackslash}p{(\linewidth - 8\tabcolsep) * \real{0.4061} + 2\tabcolsep}@{}}{%
Baixo (marcação é do provedor) \textbar{} Zero \textbar{}} \\
\end{longtable}
}

\subsection{5.5. Síntese Comparativa}\label{suxedntese-comparativa}

Cada modelo reflete escolhas políticas e históricas distintas. O Brasil
optou pelo modelo de autodeclaração com custo concentrado no pesquisador
--- a pior combinação possível: alto custo de conformidade para quem
cumpre, zero custo para quem descumpre, e impossibilidade técnica de
verificação.

A China optou pelo modelo de padrão técnico obrigatório na fonte ---
mais eficaz, mas dependente de controle centralizado sobre provedores de
IA e compatível com regime de vigilância digital. A UE optou pelo modelo
de regulação híbrida, impondo obrigações sobre provedores (marcação) e
sobre usuários (transparência em certos usos), com abordagem baseada em
risco. Os EUA optaram pelo modelo de não-regulação específica, tratando
IA como ferramenta de pesquisa e focando segurança em vez de integridade
textual.

A Portaria CNPq nº 2.664/2026, neste cenário, é o único instrumento que
impõe obrigação sem oferecer meio de cumprimento. É regulação sem
tecnologia.

\section{6. Conclusão: Balanço Crítico e
Propostas}\label{conclusuxe3o-balanuxe7o-cruxedtico-e-propostas}

\subsection{6.1. Acertos da Portaria}\label{acertos-da-portaria}

Não obstante as críticas desenvolvidas ao longo desta dissertação, a
Portaria apresenta méritos que devem ser reconhecidos. Em primeiro
lugar, o reconhecimento institucional: a Portaria tira a IAG da
invisibilidade normativa e a reconhece como fenômeno relevante na
pesquisa --- mérito não trivial em um cenário de rápida transformação
tecnológica. Em segundo lugar, a opção pela não proibição:
diferentemente de posições reacionárias que propugnam pela vedação total
do uso de IA, a Portaria busca regular o uso responsável. Em terceiro
lugar, a atualização do debate: a Portaria acelera a discussão sobre
integridade em universidades e periódicos. Em quarto lugar, o
alinhamento internacional: ainda que imperfeita, a Portaria se insere em
movimento global de regulação da IA na ciência.

\subsection{6.2. Fragilidades}\label{fragilidades}

As fragilidades identificadas são, contudo, substantivas e estruturais.
A inaplicabilidade técnica dos detectores de IAG, demonstrada por
Weber-Wulff et al.~(2023), Liu et al.~(2024), Pudasaini et al.~(2025) e
Liang et al.~(2025), invalida a premissa central da fiscalização. A
proporção uso/declaração de 40:1 (Liang et al., 2025) é o dado mais
contundente contra a eficácia do protocolo.

O custo burocrático imposto ao pesquisador não encontra contrapartida em
ganho demonstrável de integridade científica. A assimetria punitiva ---
que penaliza quem declara e não quem omite --- produz seleção adversa.
As sete ambiguidades do art. 9º geram insegurança jurídica. As colisões
normativas com a LGPD e a Constituição Federal não são resolvidas. A
vedação a pareceres com IAG (alínea ``e'') é facilmente contornada por
modelos locais. E, diferentemente de China (GB 45438-2025) e UE (Art. 50
AI Act), o Brasil não exige marcação técnica pelos provedores.

\subsection{6.3. Propostas de
Aperfeiçoamento}\label{propostas-de-aperfeiuxe7oamento}

Com base na análise desenvolvida, propõem-se as seguintes recomendações:

\begin{enumerate}
\def\labelenumi{\arabic{enumi}.}
\item
  Adoção de modelo híbrido nos moldes da UE: obrigação de marcação
  técnica pelos provedores de IA (LLMs comerciais) somada à declaração
  simplificada pelo pesquisador, transferindo o ônus da verificação para
  quem tem capacidade técnica de implementá-la.
\item
  Graduação das exigências por nível de risco, distinguindo assistência
  linguística (nível 1), sumarização e tradução (nível 2), geração
  substantiva de seções (nível 3) e geração integral (nível 4), com
  obrigações proporcionais ao risco.
\item
  Substituição da declaração detalhada por checklist padronizado de uma
  página, nos moldes adotados por periódicos de alto impacto, reduzindo
  o custo de conformidade.
\item
  Estabelecimento de sanções claras e proporcionais, respeitando o
  princípio da legalidade estrita, com tipificação objetiva de infrações
  e gradientes de penalidade.
\item
  Harmonização expressa com a LGPD e o art. 218 da CF, incluindo
  orientação específica sobre proteção de dados e proporcionalidade
  regulatória.
\item
  Negociação com provedores de IA (OpenAI, Google, Meta, DeepSeek,
  Mistral) para implementação de marcação automática de conteúdo gerado
  em português, aproveitando a capilaridade do mercado lusófono.
\end{enumerate}

\subsection{6.4. Considerações Finais}\label{considerauxe7uxf5es-finais}

Esta dissertação demonstrou que a Portaria CNPq nº 2.664/2026, não
obstante suas boas intenções, opera como instrumento de retardamento
burocrático da pesquisa científica brasileira. As evidências são
convergentes: detectores falham, a proporção uso/declaração é de 40:1, o
texto contém ambiguidades insanáveis, a fiscalização é inviável, e o
modelo brasileiro é o que oferece menor capacidade de enforcement entre
as quatro principais jurisdições analisadas.

O paradoxo central é que a Portaria tenta resolver um problema técnico
--- a detecção de IA --- com instrumento jurídico --- a declaração
obrigatória. Este descompasso entre a natureza do problema e o meio de
solução condena o protocolo à ineficácia. O pesquisador que declarar o
uso de IA será escrutinado; o que não declarar, não. O inovador honesto
paga o custo burocrático; o mal-intencionado paga zero. A Portaria não
protege a integridade: apenas a burocratiza.

A saída não é abandonar a integridade, mas redefini-la. Avaliação
substantiva baseada em impacto replicável e relevância social será
sempre mais eficaz que declarações inverificáveis. Enquanto o CNPq não
enfrentar o problema real --- a inviabilidade técnica da detecção ---, a
ciência brasileira continuará pagando o preço de um protocolo que não
protege, mas apenas retarda.

\section{Referências}\label{referuxeancias}

ALBUQUERQUE, U. P. Ciência em debate: o peso silencioso da avaliação de
cursos e publicações científicas. \textbf{The Conversation}, 13 out.
2025. Disponível em:
\url{https://theconversation.com/ciencia-em-debate-o-peso-silencioso-da-avaliacao-de-cursos-e-publicacoes-cientificas-266875}

BELADIYA, R. How China Is Controlling Generative AI Technologies.
\textbf{AI Law Guide}, maio 2026. Disponível em:
\url{https://ailawguide.org/blog/how-china-is-controlling-generative-ai-technologies}

BRASIL. Constituição da República Federativa do Brasil de 1988. Art. 5º,
II e XXXIX; Art. 218.

BRASIL. Lei nº 13.709, de 14 de agosto de 2018 (Lei Geral de Proteção de
Dados Pessoais).

BRASIL. Portaria CNPq nº 2.664, de 6 de março de 2026. Institui a
Política de Integridade na Atividade Científica. \textbf{Diário Oficial
da União}, 11 mar. 2026. Disponível em:
\url{http://memoria2.cnpq.br/web/guest/view/-/journal_content/56_INSTANCE_0oED/10157/23142775}

BRASIL, A. L. \textbf{A inteligência artificial na pesquisa e no
fomento}: desafios e oportunidades. Brasília: CAPES, 2026.

BUCCI, M. P. D. \textbf{Fundamentos para uma teoria jurídica das
políticas públicas}. 2. ed.~São Paulo: Saraiva Jur, 2021.

CABALLERO RIVERO, A.; SANTOS, R. N. M.; TRZESNIAK, P. Efeitos dos
sistemas de avaliação de pesquisa de CAPES e CNPQ nos padrões de
publicação dos pesquisadores das ciências da saúde no Brasil. \textbf{Em
Questão}, Porto Alegre, v. 30, 2024. DOI: 10.1590/1808-5245.30.138437

CAPES. Ofício Circular nº 46/2024. Diretrizes para o ciclo avaliativo
2025-2028. Brasília, 2024.

CHINA. Ministry of Science and Technology. Guidelines for Responsible
Research Conduct. Beijing, 2023. Disponível em:
\url{http://english.www.gov.cn/news/202401/06/content_WS6598c927c6d0868f4e8e2d24.html}

CHINA. Measures for the Identification of AI-Generated (Synthetic)
Content. CAC/MIIT/MPS/NRTA, 2025. Padrão GB 45438-2025. Vigente desde
01/09/2025.

EUROPEAN UNION. Regulation (EU) 2024/1689 (Artificial Intelligence Act).
\textbf{Official Journal of the European Union}, 2024. Disponível em:
\url{https://eur-lex.europa.eu/eli/reg/2024/1689}

GILS, T. et al.~From Policy to Practice: Prototyping The EU AI
Act\textquotesingle s Transparency Requirements. \textbf{SSRN}, 2024.
DOI: 10.2139/ssrn.4714345

JUSTEN FILHO, M. \textbf{Curso de direito administrativo}. 15. ed.~Rio
de Janeiro: Forense, 2024.

KOCAK, B. et al.~Ensuring peer review integrity in the era of large
language models. \textbf{European Journal of Radiology Artificial
Intelligence}, v. 2, 100018, 2025. DOI: 10.1016/j.ejrai.2025.100018

KWON, D. Science sleuths flag hundreds of papers that use AI without
disclosing it. \textbf{Nature}, v. 641, p.~290-291, 2025. DOI:
10.1038/d41586-025-01180-2

LIANG, W. et al.~Discrepancy between AI content proportion and
disclosure rate in scientific publications. \textbf{Nature Human
Behaviour}, 2025. DOI: 10.48550/arXiv.2512.06705

LIU, J. Q. J. et al.~The great detectives: humans versus AI detectors in
catching large language model-generated medical writing.
\textbf{International Journal for Educational Integrity}, v. 20, n.~1,
2024. DOI: 10.1007/s40979-024-00155-6

MESZAROS, J.; HUYS, I.; IOANNIDIS, J. P. A. Challenges in applying the
EU AI Act research exemptions to contemporary AI research. \textbf{npj
Digital Medicine}, 2026. DOI: 10.1038/s41746-025-02263-0

NSF. NSF Scientific Integrity Policy (NSF 24-007). 2024. Disponível em:
\url{https://nsf-gov-resources.nsf.gov/pubs/2024/nsf24007/nsf24007.pdf}

NSF. Trusted Research Using Safeguards and Transparency (TRUST). 2024.
Disponível em:
\url{https://nsf-gov-resources.nsf.gov/files/NSF\%20OCRSSP\%20TRUST\%20Policy\%20Memo.pdf}

PEREZ, O. C. Avaliação em Disputa: A Reforma do Qualis e os Desafios
para a Ciência. \textbf{Novos Debates}, v. 11, n.~1, 2025. DOI:
10.48006/2358-0097/V11N1.E111011

PÉREZ-NERI, I. et al.~Detecting and correcting errors in scientific
literature in the generative AI era. \textbf{Frontiers in Artificial
Intelligence}, v. 8, 2025. DOI: 10.3389/frai.2025.1644098

PESQUISAI. CNPq reconhece o uso da Inteligência Artificial Generativa na
pesquisa científica. \textbf{Portal PesquisAI}, 16 mar. 2026. Disponível
em:
\url{https://www.pesquisai.info/post/cnpq-reconhece-o-uso-da-inteligencia-artificial-generativa-na-pesquisa-cientifica-notas-sobre-a-por}

PUDASAINI, S. et al.~Survey on AI-generated plagiarism detection.
\textbf{Journal of Academic Ethics}, v. 23, p.~1137-1170, 2025. DOI:
10.1007/s10805-024-09576-x

SCHILKE, O.; REIMANN, M. The transparency dilemma: How AI disclosure
erodes trust. \textbf{Organizational Behavior and Human Decision
Processes}, v. 188, 2025. DOI: 10.1016/j.obhdp.2025.104405

SCHWAB, K. \textbf{A Quarta Revolução Industrial}. São Paulo: Edipro,
2016.

SGUISSARDI, V. A avaliação defensiva no ``modelo CAPES de avaliação''.
\textbf{Perspectiva}, Florianópolis, v. 24, n.~1, p.~49-88, 2006.
Disponível em: \url{https://repositorio.ufsc.br/handle/123456789/186529}

WEBER-WULFF, D. et al.~Testing of detection tools for AI-generated text.
\textbf{International Journal for Educational Integrity}, v. 19, n.~1,
2023. DOI: 10.1007/s40979-023-00146-z

\end{document}
